% !Mode:: "TeX:UTF-8"
\chapter{实验配置与运行命令}

本附录记录了本文实验使用的典型运行命令模板和关键配置。

\section{正式实验运行命令}

\begin{verbatim}
# Zero-shot
python scripts/run_rtllm_experiment.py \
  --model gpt-5.4 --mode zero-shot \
  --problems STUDY_12

# Feedback
python scripts/run_rtllm_experiment.py \
  --model gpt-5.4 --mode feedback \
  --k 3 --max-iterations 5 \
  --feedback-level succinct \
  --problems STUDY_12
\end{verbatim}

\section{环境配置}

实验环境：Python 3.11，Icarus Verilog 12.0，Ubuntu 22.04。API通过第三方统一网关接入。

\section{核心代码模块索引}

\begin{table}[htbp]
\centering
\caption{系统核心代码模块索引}
\begin{tabular}{ll}
\toprule
模块 & 文件路径 \\
\midrule
反馈循环控制器 & \texttt{src/feedback/loop\_runner.py} \\
提示词构建器 & \texttt{src/feedback/prompt\_builder.py} \\
Verilog执行器 & \texttt{src/runner/verilog\_executor.py} \\
评分器 & \texttt{src/ranking/ranker.py} \\
RTLLM加载器 & \texttt{src/runner/rtllm\_loader.py} \\
VerilogEval加载器 & \texttt{src/runner/verilogeval\_loader.py} \\
LLM客户端 & \texttt{src/llm/client.py} \\
Verilog提取器 & \texttt{src/utils/extract\_verilog.py} \\
实验产物管理 & \texttt{src/utils/artifacts.py} \\
代码审计 & \texttt{scripts/audit\_project.py} \\
数据审计 & \texttt{scripts/audit\_data.py} \\
\bottomrule
\end{tabular}
\end{table}
